\documentclass[11pt, a4paper]{article}

\usepackage[utf8]{inputenc}
\usepackage[T1]{fontenc}
\usepackage[english]{babel}
\usepackage{geometry}
\usepackage{graphicx}
\usepackage{booktabs}
\usepackage{array}
\usepackage{parskip}
\usepackage{hyperref}
\usepackage{microtype}
\usepackage{lmodern}

\geometry{
    top=2.5cm,
    bottom=2.5cm,
    left=2.5cm,
    right=2.5cm
}

\hypersetup{
    colorlinks=true,
    linkcolor=black,
    urlcolor=blue,
    citecolor=black
}

% -------------------------------------------------------
\title{\textbf{Bag-of-words classifier}\\[0.4em]
       \large Project 2 --- Computer Vision and Pattern Recognition}
\author{Elisabetta Chisso (IN2000231)}
\date{Summer 2026}
% -------------------------------------------------------

\begin{document}

\maketitle
\thispagestyle{empty}

% -------------------------------------------------------
\section{Introduction}

This project implements an image classifier based on the Bag of
Words (BoW) approach, applied to the 15-scene dataset by Lazebnik et
al.~(2006). The dataset contains 1,500 training images and 2,985 test
images across 15 indoor and outdoor scene categories.

The goal is to build a complete pipeline that extracts local SIFT
descriptors from images, builds a visual vocabulary using k-means
clustering, and represents each image as a normalised histogram of
visual words. Two classifiers are compared: a 1-NN baseline and a
linear one-vs-rest SVM.

% -------------------------------------------------------
\section{Pipeline}

The pipeline consists of five main steps:
\begin{itemize}
    \item Dataset loading with consistent labels across training and test splits.
    \item SIFT descriptor extraction from each image.
    \item Visual vocabulary construction via k-means on training descriptors.
    \item Image encoding as an L1-normalised histogram of visual words.
    \item Classification with 1-NN (baseline) and linear one-vs-rest SVM.
\end{itemize}

\subsection{SIFT descriptor extraction}

For each image, OpenCV's SIFT detector automatically identifies
keypoints and computes a 128-dimensional descriptor for each one. SIFT
descriptors are invariant to scale and rotation, which is important for
a heterogeneous dataset like 15-scene.

From the training set (1,275 images after the validation split),
approximately 640,000 descriptors are extracted in total, averaging
around 500 per image.

\subsection{Visual vocabulary}

All descriptors extracted from the training images are concatenated and
clustered using \texttt{MiniBatchKMeans}. Mini-batch k-means was
preferred over standard k-means for scalability: with approximately
640,000 descriptors, standard k-means iterations over the full dataset
would be prohibitively slow, whereas mini-batch updates provide a
close approximation at a fraction of the computational cost. The
resulting $k$ centroids form the visual vocabulary: the ``visual
words'' that represent recurring local patterns across scenes.

The choice of $k$ involves a trade-off: a vocabulary that is too small
produces a coarse representation, while one that is too large yields
sparse histograms that are harder to compare. For this reason, $k$ is
selected using a validation set (see Section~\ref{sec:selection}).

\subsection{Encoding and normalisation}

Each image is represented as a histogram of $k$ bins: every descriptor
is assigned to its nearest visual word (closest centroid), and the
number of occurrences of each word is counted. The histogram is then
L1-normalised, converting raw counts into relative frequencies that are
independent of the number of descriptors in the image.

\subsection{Classifiers}

\textbf{1-NN (baseline):} a 1-nearest-neighbour classifier using
Euclidean distance. Each test image is assigned to the category of the
training image whose histogram is most similar. It requires no training
phase but is sensitive to noise in the histograms.

\textbf{Linear SVM one-vs-rest:} a binary classifier is trained for
each category. To classify a test image, the real-valued outputs of all
15 classifiers are computed and the image is assigned to the class with
the highest score. The SVM is more robust because it learns a global
decision boundary from all training examples. The regularisation
parameter $C$ was left at its default value of 1.0 and was not
additionally tuned, as the focus of model selection was on the
vocabulary size $k$.

% -------------------------------------------------------
\section{Vocabulary size selection}
\label{sec:selection}

15\% of the training set was held out as a validation set (225 images)
before descriptor extraction, using stratified sampling to preserve
class balance. The vocabulary is built exclusively from the remaining
1,275 training images, and the test set is never used during model
selection.

Table~\ref{tab:val} reports validation accuracy for four values of $k$.

\begin{table}[h]
\centering
\caption{Validation accuracy for different vocabulary sizes.}
\label{tab:val}
\begin{tabular}{cccc}
\toprule
$k$ & 1-NN (val) & SVM (val) & Note \\
\midrule
50  & 32.9\% & 36.0\% & \\
\textbf{100} & 32.9\% & \textbf{42.2\%} & best \\
200 & 30.7\% & 40.9\% & \\
500 & 28.0\% & 40.0\% & \\
\bottomrule
\end{tabular}
\end{table}

Figure~\ref{fig:accuracy} shows validation accuracy as a function of
vocabulary size for both classifiers.

\begin{figure}[h]
    \centering
    \includegraphics[width=0.75\textwidth]{accuracy_vs_k.png}
    \caption{Validation accuracy vs vocabulary size $k$.}
    \label{fig:accuracy}
\end{figure}

$k{=}100$ achieves the best SVM validation accuracy (42.2\%) and is
therefore selected for the final evaluation. Notably, 1-NN does not
benefit from larger vocabularies: as $k$ increases, 1-NN accuracy
drops, likely because histograms become sparser and Euclidean distance
loses discriminative power. The SVM is more stable across vocabulary
sizes, thanks to its regularisation mechanism.

% -------------------------------------------------------
\section{Results on the test set}

With $k{=}100$, both classifiers were retrained on the full training
set (1,500 images, combining the training and validation subsets) and
evaluated on the test set (2,985 images).
Results are reported in Table~\ref{tab:test}.

\begin{table}[h]
\centering
\caption{Test set accuracy with the selected vocabulary size $k{=}100$.}
\label{tab:test}
\begin{tabular}{lc}
\toprule
Classifier & Test accuracy \\
\midrule
1-NN              & 35.0\% \\
\textbf{Linear SVM} & \textbf{42.9\%} \\
\bottomrule
\end{tabular}
\end{table}

Figures~\ref{fig:cm_knn} and~\ref{fig:cm_svm} show the normalised
confusion matrices on the test set for the 1-NN and SVM classifiers
respectively, with $k{=}100$.

\begin{figure}[h]
    \centering
    \includegraphics[width=0.82\textwidth]{confusion_matrix_1-NN_k100.png}
    \caption{Confusion matrix --- 1-NN, $k{=}100$ (test set).}
    \label{fig:cm_knn}
\end{figure}

\begin{figure}[h]
    \centering
    \includegraphics[width=0.82\textwidth]{confusion_matrix_Linear_SVM_k100.png}
    \caption{Confusion matrix --- Linear SVM, $k{=}100$ (test set).}
    \label{fig:cm_svm}
\end{figure}

The SVM outperforms the 1-NN baseline by approximately 8 percentage
points. This result is consistent with expectations for a standard BoW
pipeline on a moderately difficult dataset (a random classifier would
achieve around 6.7\% on 15 classes).

Inspecting the confusion matrices, the most commonly confused
categories are those that share similar low-level patch textures. Open
outdoor scenes (Coast, OpenCountry, Mountain) tend to be confused with
each other because they contain large uniform regions (sky, low
vegetation, ground). Similarly, indoor categories (Office, Kitchen,
Store) are confused due to rectangular objects and regular textures.
These errors are expected in a BoW system that discards all spatial
information.

% -------------------------------------------------------
\section{Conclusions}

The pipeline is intentionally minimal and transparent: no descriptor
caching, no advanced normalisation, no off-the-shelf BoW library. Each
step is implemented with small, well-documented functions that can be
read and verified line by line.

The final result (42.9\% with linear SVM, $k{=}100$) is in line with
expectations for a standard BoW baseline. The main limitations are the
lack of spatial information and the loss of structure during the
pooling step (the histogram does not encode where in the image each
visual word appears). Natural extensions would include Spatial Pyramid
Matching~(Lazebnik et al., 2006) or a chi-squared kernel for the SVM.

% -------------------------------------------------------
\appendix
\section*{Appendix: AI usage}
\addcontentsline{toc}{section}{Appendix: AI usage}

Claude was used as a support tool during the development of
this project. Specifically:
\begin{itemize}
    \item \textbf{Code structure:} Claude suggested organising the code
    into separate modules (\texttt{features.py}, \texttt{vocabulary.py},
    \texttt{classifier.py}, \texttt{utils.py}) and helped identify and
    fix several bugs in the pipeline (non-deterministic category
    ordering, descriptor/label misalignment after the train/validation
    split, figure save paths).
    \item \textbf{Comments and documentation:} Claude helped refine code
    comments and notebook Markdown cells to make them clearer and more
    readable for someone unfamiliar with the project.
\end{itemize}

All implementation choices and experimental results were independently
verified and understood.

\end{document}